% tools/build_docs_pdf.py --engine latex 用的版式。
% 文档类 ctexbook(中文断行/标点/字体由 ctex 管),oneside,A4。
% 想调样式只改这个文件:颜色、字体、标题、引用块、代码块、表格各占一段。

\usepackage{xcolor}
\usepackage{etoolbox}
\usepackage{newunicodechar}
\usepackage{enumitem}
\usepackage{fancyhdr}
\usepackage{titlesec}
\usepackage{mdframed}
\usepackage{fvextra}
\usepackage{xurl}
\usepackage{booktabs}
\usepackage{longtable}
\usepackage{array}
\usepackage{float}
\usepackage{caption}

% ---------------------------------------------------------------- 颜色
\definecolor{ink}{HTML}{1F2328}
\definecolor{muted}{HTML}{59636E}
\definecolor{linkblue}{HTML}{0969DA}
\definecolor{ruleline}{HTML}{D0D7DE}
\definecolor{softbg}{HTML}{F6F8FA}
\definecolor{quotebar}{HTML}{AFB8C1}
\definecolor{accent}{HTML}{2C6FBB}
\definecolor{warncol}{HTML}{BF8700}
\definecolor{okcol}{HTML}{1A7F37}
\definecolor{nocol}{HTML}{CF222E}

% ---------------------------------------------------------------- 字体
% 正文宋体(Songti SC 有粗体),标题无衬线(苹方有 Semibold),代码等宽。
\setCJKmainfont{Songti SC}[BoldFont={Songti SC Bold}, ItalicFont={Kaiti SC}]
\setCJKsansfont{PingFang SC}[BoldFont={PingFang SC Semibold}]
\setCJKmonofont{Songti SC}
\setmainfont{STIX Two Text}
\setsansfont{Helvetica Neue}
\setmonofont{Menlo}[Scale=0.86]
\newfontfamily\symfont{Apple Symbols}
% 正文字体(STIX Two Text)和代码字体(Menlo)都没有的符号:箭头、数学运算符、带圈数字。
% 箭头/运算符借 STIX Two Math 的字形;①~⑳ 交给 CJK 字体(宋体里有)。
\newfontfamily\mathsymfont{STIX Two Math}
\newcommand{\usym}[1]{\ifmmode\text{{\mathsymfont #1}}\else{\mathsymfont #1}\fi}
\newunicodechar{‖}{\usym{‖}}
\newunicodechar{⁄}{\usym{⁄}}
\newunicodechar{←}{\usym{←}}
\newunicodechar{↑}{\usym{↑}}
\newunicodechar{→}{\usym{→}}
\newunicodechar{↓}{\usym{↓}}
\newunicodechar{↔}{\usym{↔}}
\newunicodechar{⇔}{\usym{⇔}}
\newunicodechar{⇒}{\usym{⇒}}
\newunicodechar{⇐}{\usym{⇐}}
\newunicodechar{∂}{\usym{∂}}
\newunicodechar{∇}{\usym{∇}}
\newunicodechar{√}{\usym{√}}
\newunicodechar{∞}{\usym{∞}}
\newunicodechar{≈}{\usym{≈}}
\newunicodechar{≠}{\usym{≠}}
\newunicodechar{≡}{\usym{≡}}
\newunicodechar{≤}{\usym{≤}}
\newunicodechar{≥}{\usym{≥}}
\newunicodechar{⬜}{\usym{⬜}}
\xeCJKDeclareCharClass{CJK}{"222B}   % ∫:交给宋体(newunicodechar 会和 unicode-math 打架)
\newfontfamily\cjksymfont{Songti SC}
\newunicodechar{①}{{\cjksymfont ①}}
\newunicodechar{②}{{\cjksymfont ②}}
\newunicodechar{③}{{\cjksymfont ③}}
\newunicodechar{④}{{\cjksymfont ④}}
\newunicodechar{⑤}{{\cjksymfont ⑤}}
\newunicodechar{⑥}{{\cjksymfont ⑥}}
\newunicodechar{⑦}{{\cjksymfont ⑦}}
\newunicodechar{⑧}{{\cjksymfont ⑧}}
\newunicodechar{⑨}{{\cjksymfont ⑨}}
\xeCJKDeclareCharClass{CJK}{"258D}
\newfontfamily\dingfont{Zapf Dingbats}
% 笔记里的 ⚠️ ✅ ❌:彩色 emoji 排不出来,换成同色的排版符号
\newcommand{\wmark}{{\symfont\color{warncol}⚠}\,}
\newcommand{\okmark}{{\dingfont\color{okcol}✓}}
\newcommand{\nomark}{{\dingfont\color{nocol}✗}}
\newunicodechar{✅}{\okmark}
\newunicodechar{✓}{\okmark}
\newunicodechar{✗}{\nomark}
\newunicodechar{❌}{\nomark}

% ---------------------------------------------------------------- 版心与段落
\linespread{1.28}
\setlength{\emergencystretch}{3em}
\raggedbottom
\setlist{itemsep=0.2em, topsep=0.3em, parsep=0pt}
\floatplacement{figure}{H}
% 图片最高占版心的 72%:pandoc 默认缩到整个版心高,再加图注就装不下;
% [H] 的图又不许浮走,xelatex 会一页接一页地空转。这里把 pandoc 的 \pandocbounded 改成 0.72。
% 整本书一张图都没有时 pandoc 不定义它(也不加载 graphicx),此时跳过。
\makeatletter
\@ifundefined{pandocbounded}{}{%
\renewcommand*\pandocbounded[1]{%
  \sbox\pandoc@box{#1}%
  \Gscale@div\@tempa{0.72\textheight}{\dimexpr\ht\pandoc@box+\dp\pandoc@box\relax}%
  \Gscale@div\@tempb{\linewidth}{\wd\pandoc@box}%
  \ifdim\@tempb\p@<\@tempa\p@\let\@tempa\@tempb\fi
  \ifdim\@tempa\p@<\p@\scalebox{\@tempa}{\usebox\pandoc@box}%
  \else\usebox{\pandoc@box}\fi}}
\makeatother
\captionsetup{font={small,color=muted}, labelfont={small,color=muted}, labelsep=quad, skip=6pt}

% ---------------------------------------------------------------- 标题
% part = 分组扉页;chapter = 一篇笔记;section ~ paragraph 对应笔记里的 ## ~ ####
\titleformat{\part}[display]{\centering\sffamily\bfseries}{}{0pt}{\fontsize{26}{36}\selectfont\color{ink}}[]
\titlespacing*{\part}{0pt}{0.30\textheight}{0pt}
\titleformat{\chapter}[block]{\sffamily\bfseries\fontsize{20}{28}\selectfont\color{ink}}{}{0pt}{}[{\vspace{3pt}\color{ruleline}\titlerule[1.2pt]}]
\titlespacing*{\chapter}{0pt}{-6pt}{16pt}
\titleformat{\section}{\sffamily\bfseries\Large\color{ink}}{}{0pt}{\color{accent}\rule[-0.22em]{3pt}{1.25em}\hspace{0.55em}}
\titlespacing*{\section}{0pt}{2.2ex plus 1ex minus .3ex}{1.2ex plus .2ex}
\titleformat{\subsection}{\sffamily\bfseries\large\color{ink}}{}{0pt}{}
\titlespacing*{\subsection}{0pt}{2ex plus .8ex minus .2ex}{0.9ex plus .2ex}
\titleformat{\subsubsection}{\sffamily\bfseries\normalsize\color{ink}}{}{0pt}{}
\titlespacing*{\subsubsection}{0pt}{1.6ex plus .6ex minus .2ex}{0.6ex plus .2ex}
\titleformat{\paragraph}[runin]{\sffamily\bfseries\normalsize\color{muted}}{}{0pt}{}
\titlespacing*{\paragraph}{0pt}{1.2ex}{0.8em}

% ---------------------------------------------------------------- 页眉页脚
\pagestyle{fancy}
\fancyhf{}
\renewcommand{\chaptermark}[1]{\markboth{#1}{}}
\setlength{\headheight}{15pt}
\fancyhead[L]{\small\sffamily\color{muted}\leftmark}
\fancyfoot[C]{\small\color{muted}\thepage}
\renewcommand{\headrulewidth}{0.4pt}
\renewcommand{\headrule}{\hbox to\headwidth{\color{ruleline}\leaders\hrule height \headrulewidth\hfill}}
\fancypagestyle{plain}{\fancyhf{}\fancyfoot[C]{\small\color{muted}\thepage}\renewcommand{\headrulewidth}{0pt}}

% ---------------------------------------------------------------- 引用块:左边一条竖线 + 浅底
\renewenvironment{quote}{%
  \begin{mdframed}[hidealllines=true, leftline=true, linewidth=2.6pt, linecolor=quotebar,
    backgroundcolor=softbg, innerleftmargin=9pt, innerrightmargin=8pt,
    innertopmargin=4pt, innerbottommargin=4pt, skipabove=7pt, skipbelow=7pt]%
  \color{ink}\small
}{\end{mdframed}}

% 引用块里装了表格时(「怎么读这篇」):不可断页,见 pdf_latex_filter.lua
\newenvironment{tablequote}{%
  \par\vspace{6pt}\noindent\begin{minipage}{\linewidth}%
  \begin{mdframed}[hidealllines=true, leftline=true, linewidth=2.6pt, linecolor=quotebar,
    backgroundcolor=softbg, innerleftmargin=9pt, innerrightmargin=8pt,
    innertopmargin=4pt, innerbottommargin=4pt, nobreak=true]%
  \color{ink}\small
}{\end{mdframed}\end{minipage}\par\vspace{6pt}}

% ---------------------------------------------------------------- 代码块
% ```text 这类 pandoc 当普通 verbatim 处理的块:换成会自动折行、带左边线的版本
\fvset{fontsize=\small}
\RecustomVerbatimEnvironment{verbatim}{Verbatim}{breaklines, breakanywhere, frame=leftline,
  framerule=1.6pt, rulecolor=\color{quotebar}, framesep=7pt, formatcom=\color{ink}}
% 带语法高亮的块(Highlighting)同样折行;Shaded 的底色由 pandoc 自带的 framed 负责
\fvset{breaklines, breakanywhere}
\definecolor{shadecolor}{HTML}{F6F8FA}

% ---------------------------------------------------------------- 表格
\renewcommand{\arraystretch}{1.32}
\setlength{\tabcolsep}{5.5pt}
\setlength{\heavyrulewidth}{0.9pt}
\setlength{\lightrulewidth}{0.45pt}
\AtBeginEnvironment{longtable}{\small\color{ink}}
\setlength{\LTpre}{8pt}
\setlength{\LTpost}{8pt}

% ---------------------------------------------------------------- PDF 书签
% 目录页只展开到「篇」(tocdepth=0),但书签栏要全量(节、小节都有),方便跳转
\AtBeginDocument{\bookmarksetup{depth=5, open, openlevel=2}}
